% =====================================================================
%  capitulos/cap09.tex — Capítulo 9. Los predicativos
%  Autores: Pablo Zdrojewski y Fernando Carranza
%  Fuente: Capitulo_09_normalizado.docx
% =====================================================================
\capitulo{Los predicativos}{Pablo Zdrojewski y Fernando Carranza}

\section{Introducción}

El \termino{predicativo} es uno de los modificadores del núcleo verbal más
difíciles de caracterizar porque, en primer lugar, no está restringido
categorialmente, es decir, su núcleo puede pertenecer a más de una clase de
palabras y, en consecuencia, puede formar distintos tipos de sintagmas: SN, SA,
SP e, incluso, SAdv.

\begin{exe}
  \ex\label{ex:09-cat1}
    \begin{xlist}
      \ex\label{ex:09-cat1-a} Pedro es [\etq{SN} un arquitecto muy renombrado].
      \ex\label{ex:09-cat1-b} Varias torres de cristal se alzan [\etq{SA}
          gallardas] en Buenos Aires.
      \ex La abogada ofició [\etq{SP} de defensora].
      \ex La cosa es [\etq{SAdv} así].
    \end{xlist}
  \ex\label{ex:09-cat2}
    \begin{xlist}
      \ex\label{ex:09-cat2-a} Denominaron a este espécimen [\etq{SN}
          \textit{Patagosaurus fariasi}].
      \ex\label{ex:09-cat2-b} Guardé la ropa [\etq{SA} muy mojada].
      \ex Que no te traten [\etq{SP} de tonto].
      \ex A Juan lo vi [\etq{SAdv} bien].
    \end{xlist}
\end{exe}

Además, tal como iremos viendo, la relación que establece el predicativo con el
núcleo del predicado varía según el tipo de verbo de que se trate. Esto lleva a
que también cambie el tipo de contribución que hace el predicativo a la sintaxis
y a la semántica de las oraciones en las que aparece. En este capítulo vamos a
abordar los distintos tipos de predicativos que se reconocen en el español. En
primer lugar, partimos de una caracterización en relación con sus propiedades
sintácticas y semánticas e introducimos las nociones de \termino{predicación
primaria} y \termino{predicación secundaria}. Luego, en segundo lugar,
distinguimos las clases de predicativos en función de su orientación hacia el
sujeto o el objeto de la oración, y trazamos la diferencia entre obligatorios y
no obligatorios. Por último, señalamos cómo se marcan sintácticamente y
presentamos casos de homonimia estructural a que dan lugar.

\section{Los predicativos y la predicación}\label{sec:09-predicacion}

\subsection{Predicación semántica y predicación sintáctica}

La denominación \forma{predicativo} alude a que se trata de un constituyente que
expresa un tipo particular de predicación.

\begin{destacado}
Es importante tener en cuenta que, aunque están en el predicado, ninguno de los
tipos de predicativo que veremos en este capítulo se considera un argumento del
verbo; en consecuencia, tampoco les corresponde ningún papel temático, ya que
integran con el verbo la predicación principal o forman una predicación
secundaria.
\end{destacado}

Según ya se ha planteado en el \capref{7}, en términos semánticos, podemos
caracterizar la noción de predicación como una relación entre un predicado y sus
argumentos. Los argumentos denotan entidades del mundo extralingüístico,
mientras que los predicados atribuyen propiedades a los argumentos o denotan
relaciones entre argumentos. En este sentido, todas las clases de palabras y de
sintagmas pueden ser \termino{predicados semánticos}, como se aprecia en los
ejemplos de \exref{ex:09-clases}:

\begin{exe}
  \ex\label{ex:09-clases}
    \begin{xlist}
      \ex Juan [canta / construye casas]. \hfill [Verbo]
      \ex Juan [es cantante / arquitecto]. \hfill [Sustantivo]
      \ex Juan [es afinado / prolijo]. \hfill [Adjetivo]
      \ex Juan [está arriba / afuera / acá / bien]. \hfill [Adverbio]
      \ex Enzo [es de River / de Córdoba]. \hfill [Preposición]
    \end{xlist}
\end{exe}

Resulta conveniente, entonces, retomar la distinción entre las nociones
semántica y sintáctica de predicado, que se introdujeron en capítulos
anteriores. El \termino{predicado semántico} es la expresión que atribuye cierto
contenido descriptivo a los argumentos. El \termino{predicado sintáctico}, por
su parte, puede definirse como el constituyente de la oración que excluye al
sujeto, que contiene al predicado semántico y cuyo núcleo típicamente expresa
las categorías de tiempo, aspecto y modo, junto con la concordancia en número y
persona con el sujeto (los constituyentes delimitados por corchetes en
\exref{ex:09-clases}) (v. \capref{6}).\pendiente{El original remite a los
corchetes de (3) para ambos predicados} En consecuencia, solo los verbos pueden
ser a la vez predicados semánticos y núcleos de predicados sintácticos, con la
excepción de ciertas estructuras sintácticamente marcadas que involucran
predicados no verbales (v. \capref{6}, §\ref{sec:06-noverbales}) ---por ejemplo,
\forma{Excelente, ese parcial / Hermosas, tus zapatillas}---
\parencite{hernanz1999suner,ditullio2014}.\pendiente{El original remite a
«cap. 2: §4»}

\begin{ejercicio}\label{ej:09-predicados}
En las siguientes oraciones resuelva lo pedido, según se muestra en el ejemplo
resuelto~1.

\begin{consignas}
  \item Identifique el predicado sintáctico entre corchetes y subraye el
        predicado semántico.
  \item Determine a qué clase de sintagma pertenece el predicado semántico que
        identificó.
\end{consignas}

\begin{exe}
  \ex Eduarda Mansilla y Lucio Mansilla eran hermanos.
  \ex Eduarda Mansilla escribió varios libros.
  \ex Eduarda Mansilla ofició de traductora.
  \ex Eduarda Mansilla fue también compositora.
  \ex Margarita Abella Caprile nació en 1901.
  \ex La poesía de Margarita Caprile era muy íntima.
  \ex Margarita Caprile fue amiga de Alfonsina Storni.
  \ex María Luisa Carnelli era de La Plata.
  \ex María Luisa Carnelli firmaba sus letras de tango con seudónimo.
\end{exe}

\begin{resuelto}
1. Eduarda Mansilla y Lucio Mansilla [eran [\etq{SN} \subrayado{hermanos}]].
\end{resuelto}
\end{ejercicio}

\subsection{Predicación primaria y predicación secundaria}\label{sec:09-primaria}

La mayor parte de los verbos del español, como \forma{vivir}, \forma{combatir} y
\forma{dar}, tienen contenido léxico, es decir, aportan significado conceptual.
Esto los habilita a funcionar sintácticamente como verbos plenos. A diferencia
de estos, algunos verbos carecen total o parcialmente de contenido conceptual y
parecen cumplir una función meramente gramatical. Estos incluyen, por un lado, a
los \termino{verbos copulativos} \forma{ser, estar, parecer}, que reciben este
nombre precisamente porque funcionan como una cópula o nexo gramatical entre un
argumento y un predicado semántico.\footnote{Para una descripción de las clases
de palabras léxicas y funcionales en un nivel inicial, ver
\textcite{giammatteo2009}, y para ampliar estos conceptos, ver
\textcite{bosque1990}.}

\begin{exe}
  \ex\label{ex:09-copulativos}
    \begin{xlist}
      \ex María \textit{es} inteligente.
      \ex Juan \textit{está} enfermo.
      \ex Pedro \textit{parece} trabajador.
    \end{xlist}
\end{exe}

Por otro lado, esta clase abarca también a los \termino{verbos semicopulativos}
(o \termino{pseudocopulativos}), como \forma{resultar, terminar,
seguir},\footnote{No debe confundirse el uso semicopulativo de estos verbos, que
es el que exige un predicativo, con el uso pleno, en el que los verbos presentan
estructura argumental, como en \forma{terminó la tarea} o \forma{sigue por el
camino amarillo}.} que, si bien no están absolutamente desprovistos de contenido
conceptual, no funcionan como verbos plenos, pues en estos usos carecen de
estructura argumental y requieren el predicativo para completar su significado.

\begin{exe}
  \ex
    \begin{xlist}
      \ex Juan \textit{resultó} herido. / \juicio{*}{Juan resultó.}
      \ex Pedro \textit{terminó} loco.
      \ex Esteban \textit{sigue} enfermo.
    \end{xlist}
\end{exe}

Al perder su contenido léxico, estos dos tipos de verbos, si bien son núcleos de
predicados sintácticos, no constituyen predicados semánticos. La única
predicación de la oración es la que expresa el predicativo que estos verbos
exigen. Se dice, entonces, que el predicativo, como \forma{inteligente},
\forma{enfermo} y \forma{trabajador} en \exref{ex:09-copulativos}, expresa una
\termino{predicación primaria}.\pendiente{El original remite a (5)}

A diferencia de los predicativos con verbos copulativos y semicopulativos,
cuando el núcleo del predicado sintáctico es un verbo con significado léxico
pleno, este núcleo aporta la predicación principal de la oración ---es decir, la
predicación primaria---, mientras que el predicativo contribuye con una
predicación adicional ---es decir, una \termino{predicación secundaria}---, como
se ve en los ejemplos de \exref{ex:09-secundaria}.

\begin{exe}
  \ex\label{ex:09-secundaria}
    \begin{xlist}
      \ex Juan \textit{habló} dormido.\\
          Predicación primaria: \forma{Juan habló}.\\
          Predicación secundaria: \forma{Juan estaba dormido}.
      \ex María \textit{lee} muy concentrada.\\
          Predicación primaria: \forma{María lee}.\\
          Predicación secundaria: \forma{María está concentrada}.
      \ex San Martín \textit{cruzó} los Andes enfermo.\\
          Predicación primaria: \forma{San Martín cruzó los Andes}.\\
          Predicación secundaria: \forma{San Martín estaba enfermo}.
    \end{xlist}
\end{exe}

\begin{ejercicio}\label{ej:09-primaria}
Identifique en estas oraciones la predicación primaria y la secundaria (si la
hubiera), como se muestra en los ejemplos resueltos de abajo.

\begin{exe}
  \ex Carlos recibió la carta azorado.
  \ex Romina quedó perpleja.
  \ex La sala está llena.
  \ex El niño lloró desconsolado.
  \ex Las Malvinas son argentinas.
  \ex Ella respondió la llamada muy tranquila.
  \ex Martín se siente solo.
\end{exe}

\begin{resuelto}[Ejemplos resueltos]
1. Predicación primaria: \forma{Carlos recibió la carta}. Predicación
secundaria: \forma{Carlos estaba azorado}. 2. Predicación primaria:
\forma{Romina quedó perpleja}. Predicación secundaria: no hay.
\end{resuelto}
\end{ejercicio}

\section{Tipos de predicativos}\label{sec:09-tipos}

Tal como mencionamos en la introducción, es posible distinguir diferentes clases
de predicativos según dos criterios: por un lado, en función de su orientación
hacia el sujeto o hacia el objeto directo de la oración y, por el otro, en
relación con que su presencia sea obligatoria o facultativa. En lo que sigue,
presentamos ambos criterios y luego mostramos cómo se integran ambas
distinciones.

\subsection{Predicativos subjetivos y objetivos}\label{sec:09-subjetivos}

Por un lado, el predicativo es un constituyente interno del predicado verbal,
por lo que está asociado al núcleo. Por otra parte, se caracteriza por
establecer una relación de predicación con el sujeto \exref{ex:09-pedro} o con
el objeto de la oración \exref{ex:09-ciencia}.

\begin{exe}
  \ex\label{ex:09-pedro} Pedro es un arquitecto muy renombrado.
  \ex\label{ex:09-ciencia} La ciencia denominó a este espécimen
      \textit{Patagosaurus fariasi}.
\end{exe}

Representamos informalmente esta relación mediante flechas (Suj: sujeto; Pred:
predicado; V: verbo; OD: objeto directo; Pvo: predicativo): en
\exref{ex:09-flechas-a} se trata de un \termino{predicativo subjetivo}, mientras
que en \exref{ex:09-flechas-b} tenemos un \termino{predicativo objetivo}
\parencite{demonte1999}:

\begin{exe}
  \ex\label{ex:09-flechas}
    \begin{xlist}
      \ex\label{ex:09-flechas-a}
          \begin{tikzpicture}[baseline=(s.base), node distance=0pt,
              every node/.style={inner sep=1pt}]
            \node (s) {[\etq{Suj} Pedro]};
            \node[right=4pt of s] (v) {[\etq{Pred} \etq{V}es};
            \node[right=3pt of v] (p) {\etq{Pvo}[un arquitecto muy renombrado]]};
            \draw[->, >=Stealth] (p.south) .. controls +(0,-.5) and +(0,-.5) .. (s.south);
          \end{tikzpicture}
          \vspace{3ex}
      \ex\label{ex:09-flechas-b}
          \begin{tikzpicture}[baseline=(s.base), node distance=0pt,
              every node/.style={inner sep=1pt, font=\small}]
            \node (s) {[\etq{Suj} La ciencia]};
            \node[right=2pt of s] (v) {[\etq{Pred} \etq{V}denominó};
            \node[right=3pt of v] (o) {\etq{OD}[a este espécimen]};
            \node[right=3pt of o] (p) {\etq{Pvo}[\textit{Patagosaurus fariasi}]]};
            \draw[->, >=Stealth] (p.south) .. controls +(0,-.5) and +(0,-.5) .. (o.south);
          \end{tikzpicture}
          \vspace{3ex}
    \end{xlist}
\end{exe}

\begin{ejercicio}\label{ej:09-tabla}
\begin{consignas}
  \item Según se muestra en el ejemplo resuelto~1, determine si los predicativos
        señalados entre corchetes en la tabla son SN, SA o SP.
  \item Identifique en las oraciones de la tabla los sujetos y los objetos
        directos, si los hubiere.
  \item ¿Qué marca morfosintáctica visibiliza la relación que se produce entre
        el sujeto o el objeto y el predicativo, cuando este es un SN, un SA o un
        SP?
\end{consignas}

{\small
\begin{tabularx}{\linewidth}{@{}r L L@{}}
  \toprule
  & a. Predicativo subjetivo & b. Predicativo objetivo \\
  \midrule
  1. & [\etq{Suj} Santos Vega] fue [\etq{SN} un payador argentino]. &
       [\etq{Suj} Los demás payadores] consideraban [\etq{COD} a Santos Vega]
       [\etq{SA} invencible]. \\
  2. & Gabino Ezeiza murió [pobre] en el barrio de Flores. & Se suele
       considerar a José Betinotti [el nexo entre los payadores y los primeros
       tangueros]. \\
  3. & Para Gabino Ezeiza, Pancho Luna ofició [de maestro]. & José Betinotti
       convirtió [en exitosa] la línea de las payadas que exaltaban el amor
       materno. \\
  4. & Arturo de Nava solía vestir [elegante]. & El Mercosur declaró
       [Patrimonio Cultural] al arte de la payada. \\
  5. & José Betinotti siempre recordaba [agradecido] a su amigo Pablo Vázquez
       en cada función. & El alcalde acusó a Juan Moreira [de mentiroso]. \\
  \bottomrule
\end{tabularx}}
\end{ejercicio}

\begin{ejercicio}\label{ej:09-ps-po}
Según se muestra en los ejemplos resueltos, identifique cuáles de los
siguientes predicativos son subjetivos (PS) y cuáles objetivos (PO).

\begin{exe}
  \ex Por las noches cantaba ensimismada.
  \ex Rosita Quiroga interpretaba las canciones muy arrabaleras.
  \ex José Betinotti se declaraba [admirador de Alfredo Palacios].
  \ex Para una grabación particular, Mariano Villar Sáenz Peña convocó [de
      guitarrista] a Rosita Quiroga.
  \ex Las sucesivas crisis económicas dejaron al Café de los Angelitos [muy
      deteriorado].
  \ex Ángel Vargas tenía un pulmón [muy comprometido].
  \ex Alfredo Palacios y Juan B.~Justo solían discutir [enfervorizados] en el
      Café de los Angelitos.
  \ex Rosita Quiroga se retiró de su carrera artística [todavía joven].
  \ex En sus últimos años de vida a Enrique Santos Discépolo se lo veía [solo].
  \ex Azucena Maizani murió [casi olvidada] en 1970.
\end{exe}

\begin{resuelto}[Ejemplos resueltos]
1. Por las noches cantaba [\etq{PS} ensimismada]. 2. Rosita Quiroga
interpretaba las canciones [\etq{PO} muy arrabaleras].
\end{resuelto}
\end{ejercicio}

\subsection{Predicativos obligatorios y no obligatorios}\label{sec:09-obligatorios}

Además de la distinción entre predicativos subjetivos y objetivos que vimos en la
sección anterior, los predicativos también se pueden distinguir respecto de su
relación con el verbo. Si retomamos las oraciones de \exref{ex:09-cat1} y
\exref{ex:09-cat2}, se puede apreciar que la omisión del predicativo en
\exref{ex:09-cat1-b} y \exref{ex:09-cat2-b} no afecta la gramaticalidad de la
oración ni la interpretación del verbo, como se aprecia en
\exref{ex:09-opcionales}. Los paréntesis indican la opcionalidad de los
predicativos en estos casos:

\begin{exe}
  \ex\label{ex:09-opcionales}
    \begin{xlist}
      \ex Varias torres de cristal se alzan (gallardas) en Buenos Aires.
      \ex Guardé la ropa (muy mojada).
    \end{xlist}
\end{exe}

En contraste, los ejemplos de \exref{ex:09-obligatorios} muestran que en
\exref{ex:09-cat1-a} y \exref{ex:09-cat2-a} no es posible omitir el predicativo:

\begin{exe}
  \ex\label{ex:09-obligatorios}
    \begin{xlist}
      \ex[*]{Pedro es. (cfr. \forma{Pedro es un arquitecto muy renombrado})}
      \ex[*]{Nombraron a este dinosaurio.\footnote{La oración resulta gramatical
          si se interpreta \forma{nombrar} como \sig{mencionar su nombre}
          (\forma{Cuando tomaron lista, nombraron a Carlos}) o como \sig{designar
          a alguien en un cargo vacante} (\forma{Nombraron un nuevo ministro de
          Hacienda}).} (cfr. \forma{Nombraron a este dinosaurio
          \textit{Patagosaurus fariasi}})}
    \end{xlist}
\end{exe}

Esto nos permite distinguir dos tipos de predicativos: los \termino{predicativos
obligatorios}, que son aquellos que no pueden omitirse, y los
\termino{predicativos no obligatorios}, que son opcionales.

\begin{ejercicio}\label{ej:09-pso-psno}
Todos los ejemplos que aparecen a continuación tienen predicativos subjetivos.
Identifíquelos e indique en cada caso si son subjetivos obligatorios (PSO) o no
(PSNO), según se muestra en los ejemplos resueltos de abajo.

\begin{exe}
  \ex Pedro estaba furioso por haber perdido las llaves.
  \ex Pedro golpeó la mesa furioso.
  \ex Ernesto es biólogo.
  \ex Carlos salió despeinado de su casa.
  \ex Mi perro corre medio rengo cuando se asusta.
  \ex Ese reloj parece nuevo.
  \ex Rocío miró la foto asombrada.
  \ex Enrique escribió su novela muy entusiasmado.
  \ex Martina resultó la ganadora del concurso.
  \ex La actriz interpretó el papel muy segura.
\end{exe}

\begin{resuelto}[Ejemplos resueltos]
1. Pedro estaba [\etq{PSO} furioso] por haber perdido las llaves. 2. Pedro
golpeó la mesa [\etq{PSNO} furioso].
\end{resuelto}
\end{ejercicio}

\begin{ejercicio}\label{ej:09-sonono}
Observe los siguientes ejemplos y resuelva lo pedido más abajo, según se muestra
en el ejemplo resuelto~1.

\begin{exe}
  \ex
    \begin{xlist}
      \ex Anda [velozmente].
      \ex Anda [cansado].
    \end{xlist}
  \ex
    \begin{xlist}
      \ex Sigue [su instinto].
      \ex Sigue [mal].
      \ex Sigue [en Madrid].
      \ex Sigue [molesto] por lo que le dije.
    \end{xlist}
  \ex
    \begin{xlist}
      \ex Se puso [un sombrero].
      \ex Se puso [rojo].
    \end{xlist}
  \ex
    \begin{xlist}
      \ex Volvió [temprano].
      \ex Volvió [enfermo].
      \ex Volvió [loco] a su amigo.
      \ex Volvió [una noche].
      \ex Se volvió [loco].
    \end{xlist}
  \ex
    \begin{xlist}
      \ex Juan se sintió [olor a humo] después de hacer el asado.
      \ex Juan se sintió [cansado].
    \end{xlist}
  \ex
    \begin{xlist}
      \ex Esteban se hizo [el desayuno].
      \ex Esteban se hizo [el tonto].
    \end{xlist}
  \ex
    \begin{xlist}
      \ex Ese tipo cayó [mal]. Por eso se rompió el tobillo.
      \ex Ese tipo me cae [simpático].
    \end{xlist}
  \ex
    \begin{xlist}
      \ex Jugó [muy concentrado].
      \ex Jugó [el partido].
    \end{xlist}
\end{exe}

\begin{consignas}[label=\roman*)]
  \item Indique cuáles de los sintagmas delimitados entre corchetes son
        predicativos y cuáles no.
  \item En los casos en que haya identificado un predicativo, indique el tipo de
        sintagma de que se trata (SN o SA) y si es obligatorio (PO) o no (PNO).
  \item En los casos en que no hay predicativo, especifique el tipo de sintagma
        y la función sintáctica del constituyente delimitado.
\end{consignas}

\begin{resuelto}
1. i) a. \forma{Anda [velozmente]}: no es predicativo. b. \forma{Anda
[cansado]}: es predicativo. ii) b. Anda [\etq{SA/PO} cansado]. iii) a. Anda
[\etq{SAdv/ACM} velozmente].
\end{resuelto}
\end{ejercicio}

\subsection{Integración de los dos tipos de predicativos}\label{sec:09-integracion}

La combinación de las distinciones presentadas en las dos secciones anteriores
se puede presentar en una tabla de doble entrada como la
tabla~\ref{tab:09-tipos}.

\begin{table}[htbp]
\centering\small
\caption{Tipos de predicativos}\label{tab:09-tipos}
\begin{tabularx}{\linewidth}{@{}>{\bfseries\raggedright\arraybackslash}p{24mm} L L@{}}
  \toprule
  & \textbf{Predicativo subjetivo} & \textbf{Predicativo objetivo} \\
  \midrule
  Predicativo obligatorio & Pedro es [un arquitecto muy renombrado]. & La
    comisión nombró al postulante [presidente de la cámara]. \\
  Predicativo no obligatorio & Varias torres de cristal se alzan [gallardas] en
    Buenos Aires. & Guardé la ropa [muy mojada]. \\
  \bottomrule
\end{tabularx}
\end{table}

Siguiendo la tabla, se reconocen cuatro tipos de predicativos:

\begin{itemize}
  \item predicativo subjetivo obligatorio $\rightarrow$ PSO;
  \item predicativo subjetivo no obligatorio $\rightarrow$ PSNO;
  \item predicativo objetivo obligatorio $\rightarrow$ POO;
  \item predicativo objetivo no obligatorio $\rightarrow$ PONO.
\end{itemize}

Cada uno de estos tipos se puede describir mediante un grupo de características
que pueden utilizarse como pruebas de reconocimiento \parencite{ditullio2014},
según mostraremos en los apartados siguientes.

\subsubsection{Predicativo subjetivo obligatorio}\label{sec:09-pso}

\begin{itemize}
  \item Puede ser un sintagma de cualquier clase: SA, SN, SAdv, SP.
  \item Ocurre solo con verbos copulativos (\forma{ser, estar, parecer}) y
        semicopulativos (\forma{volverse, resultar, semejar}).
  \item Puede ser reemplazado por el pronombre invariable \forma{lo}:
        \forma{Rosa Guerra fue periodista} $\rightarrow$ \forma{Rosa Guerra lo
        fue}. Esta prueba sirve únicamente en el caso de que el verbo sea
        copulativo, por lo que no se puede aplicar a semicopulativos.
  \item Es importante señalar que este pronombre \forma{lo} es invariable y no
        debe ser confundido con el pronombre clítico que puede funcionar como
        COD y que varía en género y número (\forma{lo, la, los, las}), como se
        aprecia en el siguiente contraste:
\end{itemize}

\begin{exe}
  \ex ¿Rosa Guerra fue periodista? Sí, \textit{lo / *la} fue.
  \ex Compró el diario / la revista $\rightarrow$ \textit{Lo / La} compró.
\end{exe}

\begin{itemize}
  \item Si el predicativo es un SA, debe concordar en género y número con el
        sujeto de la oración: \forma{Rosa Guerra era argentina; Eugenio
        Cambaceres era argentino; Ellas son argentinas}.
\end{itemize}

\subsubsection{Predicativo subjetivo no obligatorio}\label{sec:09-psno}

\begin{itemize}
  \item Ocurre con verbos de contenido léxico y constituye una predicación
        secundaria.
  \item Acepta la paráfrasis por una coordinación entre la predicación primaria,
        que es la que aporta el verbo principal, y la secundaria, que requiere
        el agregado de un verbo copulativo: \forma{Pedro salió a trabajar
        despeinado} $\rightarrow$ \forma{Pedro salió a trabajar y estaba
        despeinado}. Esta prueba sirve para distinguir predicativos de
        circunstanciales, especialmente cuando se trata de un adverbio
        adjetival,\footnote{Los adverbios adjetivales se introdujeron en el
        \capref{5}, §\ref{sec:05-sadv}.} como en \forma{Juan camina rápido}
        $\rightarrow$ \juicio{*}{\forma{Juan camina y está rápido}}.\footnote{El
        diagnóstico no siempre resulta claro con predicativos conformados por
        SSPP. Por ejemplo, \forma{Juan fue al trabajo con un maletín} admite la
        coordinación \forma{Juan fue al trabajo y Juan estaba con un maletín}.}
  \item Si el predicativo es un SA, concuerda en género y número con el sujeto
        de la oración: \forma{Juan Manuel de Rosas leyó furioso / *furiosa el
        libro de Sarmiento; Manuelita Rosas leyó *furioso / furiosa el libro de
        Sarmiento}.
\end{itemize}

\subsubsection{Predicativo objetivo obligatorio}\label{sec:09-poo}

\begin{itemize}
  \item Ocurre con verbos que seleccionan un objeto directo y un predicativo,
        como, por ejemplo, \forma{nombrar, llamar, designar, notar, encontrar},
        etc. Estos verbos tienden a ser ambiguos con una lectura en la que no se
        selecciona un predicativo. Nótese que \forma{nombrar a alguien profesor}
        no significa lo mismo que \forma{nombrar a alguien}.
  \item Al ser un constituyente diferente del COD, puede anteponerse a este, es
        decir que acepta la prueba de movimiento:\footnote{La prueba de
        movimiento y las demás pruebas de reconocimiento de constituyentes se
        presentaron en el \capref{5}, §\ref{sec:05-pruebas}.} \forma{Nombraron
        al nuevo jugador capitán del equipo} $\rightarrow$ \forma{Nombraron
        capitán del equipo al nuevo jugador}.
  \item Asimismo, si el COD es reemplazado por el pronombre clítico
        correspondiente (\forma{lo, la, los, las}), el predicativo no queda
        incluido en la sustitución:
\end{itemize}

\begin{exe}
  \ex Considero \textit{esta novela} [valiosa] $\rightarrow$ \textit{La}
      considero [valiosa].
\end{exe}

\begin{itemize}
  \item Si el predicativo es un SA, concuerda en género y número con el objeto
        de la oración.
\end{itemize}

\begin{exe}
  \ex Considero \textit{esta novela} [muy valiosa]. / Considero \textit{estas
      novelas} [muy valiosas].
\end{exe}

\subsubsection{Predicativo objetivo no obligatorio}\label{sec:09-pono}

\begin{itemize}
  \item Ocurre solo con verbos transitivos.
  \item Puede anteponerse al COD (movimiento).
\end{itemize}

\begin{exe}
  \ex Encontró \textit{la mesa} [sucia] $\rightarrow$ Encontró [sucia]
      \textit{la mesa}.
\end{exe}

\begin{itemize}
  \item Si el COD es reemplazado por el pronombre clítico correspondiente
        (\forma{lo, la, los, las}), el predicativo no queda incluido en la
        sustitución:
\end{itemize}

\begin{exe}
  \ex Saludó \textit{a Pedro} [sentado en el parque] $\rightarrow$ \textit{Lo}
      saludó [sentado en el parque].
\end{exe}

\begin{itemize}
  \item Si el predicativo es un SA, concuerda en género y número con el objeto
        de la oración.
\end{itemize}

\begin{exe}
  \ex Encontró \textit{la mesa} [sucia]. / Encontró \textit{las mesas}
      [sucias].
\end{exe}

\begin{recuadro}[Para ampliar]
Los verbos que exigen predicativo objetivo obligatorio pertenecen a las
siguientes clases: verbos de entendimiento (\forma{considerar, creer, imaginar,
suponer, juzgar}), verbos de percepción sensible o intelectual (\forma{ver, oír,
notar, descubrir, encontrar}), verbos designativos (\forma{designar, elegir,
nombrar, denominar, llamar}) y verbos volitivos (\forma{querer})
\parencite[134]{ditullio2014}.
\end{recuadro}

\begin{ejercicio}\label{ej:09-pruebas}
Reconozca qué pruebas de las planteadas en §\ref{sec:09-integracion} se están
aplicando en los siguientes ejemplos con el fin de demostrar que los elementos
subrayados en la primera oración de cada par son predicativos, según se muestra
en el ejemplo resuelto~1.\pendiente{El original remite a «§5»}

\begin{exe}
  \ex
    \begin{xlist}
      \ex María sintió el pollo \subrayado{soso}.
      \ex María sintió soso el pollo.
    \end{xlist}
  \ex
    \begin{xlist}
      \ex Gaspar parecía \subrayado{distante}.
      \ex Gaspar lo parecía.
    \end{xlist}
  \ex
    \begin{xlist}
      \ex Consiguió \subrayado{caras} las entradas.
      \ex Consiguió las entradas y las entradas estaban caras.
    \end{xlist}
  \ex
    \begin{xlist}
      \ex La tarea le había resultado \subrayado{fácil}.
      \ex Las tareas le habían resultado fáciles.
    \end{xlist}
  \ex
    \begin{xlist}
      \ex Gaspar encontró las zapatillas \subrayado{baratas}.
      \ex Gaspar encontró baratas las zapatillas.
    \end{xlist}
\end{exe}

\begin{resuelto}
1. Prueba de movimiento: muestra que el predicativo no está incluido en el COD,
es decir que se trata de un constituyente diferente del COD.
\end{resuelto}
\end{ejercicio}

\begin{ejercicio}\label{ej:09-enriquez}
Lea atentamente el siguiente fragmento adaptado de «Alguien camina sobre tu
tumba», de Mariana Enriquez (2009), y resuelva las consignas que aparecen a
continuación, según se muestra en el ejemplo resuelto.

\begin{consignas}
  \item Identifique cuáles de los constituyentes destacados no son
        predicativos. Justifique su respuesta.
  \item Indique cuáles de las expresiones destacadas son predicativos y
        clasifíquelos. Explicite los criterios utilizados.
\end{consignas}

% Fragmento literario protegido: se incorpora a mano desde el original.
\textoaincorporar{el fragmento adaptado de «Alguien camina sobre tu tumba», de
M.~Enriquez, desde «El ferry que lleva a Paul hasta su trabajo» hasta
«\ldots{}en ocasiones, violenta», con los constituyentes destacados entre
corchetes: \textit{cerca del puerto; el uniforme; mi novio; australiano; de
mecánico de bicicletas; capital de Australia Occidental; a larga distancia;
insólitamente estable, permanente y feliz; la capital de estado más aislada del
mundo; la ciudad de más de un millón de habitantes más aislada del mundo;
comprobable y cierto; una localidad muy cool, algo estrafalaria, favorita de
mochileros y newagers; con jardines llenos de cristales curativos y hadas}}

\begin{resuelto}
a) \forma{El uniforme} no es un predicativo, sino un COD ([\etq{COD} el
uniforme]), ya que admite la sustitución por un pronombre acusativo. b)
\forma{Mi novio} es un predicativo subjetivo obligatorio, ya que ocurre como
complemento de un verbo copulativo y admite la sustitución por el pronombre
invariable \forma{lo} (\forma{Paul lo es}).
\end{resuelto}
\end{ejercicio}

\begin{ejercicio}\label{ej:09-significado}
Según se muestra en el ejemplo resuelto~1, responda las siguientes consignas.

\begin{consignas}[label=\roman*)]
  \item ¿Cuáles de los siguientes ejemplos presentan predicativos?
  \item ¿Cómo afecta la presencia de predicativos al significado de estos
        verbos? Explique en qué consiste la diferencia de significado.
\end{consignas}

\begin{exe}
  \ex
    \begin{xlist}
      \ex Paul encontró el uniforme sobre la mesa.
      \ex Paul encontró el uniforme sobre la mesa muy arrugado.
    \end{xlist}
  \ex
    \begin{xlist}
      \ex La gente permaneció estupefacta.
      \ex La gente permaneció en el lugar.
      \ex La gente permaneció estupefacta en el lugar.
    \end{xlist}
  \ex
    \begin{xlist}
      \ex Perth está aislada.
      \ex Perth está en Australia.
    \end{xlist}
  \ex
    \begin{xlist}
      \ex El testigo declaró en el juicio.
      \ex El juez declaró culpable al imputado.
    \end{xlist}
\end{exe}

\begin{resuelto}
i) 1a. No presenta ningún predicativo. 1b. PO [muy arrugado]. ii) En 1a,
\forma{encontrar} significa \sig{hallar algo cuyo paradero se desconocía}; en
1b, la predicación describe el estado en que el objeto ---el uniforme--- se
encontró.
\end{resuelto}
\end{ejercicio}

\begin{ejercicio}\label{ej:09-vf}
Revise las respuestas del ejercicio~\ref{ej:09-primaria} y determine cuáles de
estas afirmaciones son verdaderas y cuáles falsas, según se muestra en el
ejemplo resuelto i.

\begin{enumerate}[label=\roman*.]
  \item Los PSO expresan predicaciones primarias. \hfill V / F
  \item Los PSO expresan predicaciones secundarias. \hfill V / F
  \item Los PSNO expresan predicaciones primarias. \hfill V / F
  \item Los PSNO expresan predicaciones secundarias. \hfill V / F
\end{enumerate}

\begin{resuelto}
i. V.
\end{resuelto}
\end{ejercicio}

\begin{ejercicio}\label{ej:09-verbos}
\begin{consignas}
  \item Determine cuáles de los verbos que aparecen en la siguiente lista
        requieren predicativos en alguna lectura y complete la tabla.
  \item Siempre que sea posible, construya un ejemplo con un predicativo, tal
        como se ilustra en el ejemplo resuelto~1.
\end{consignas}

\begin{center}\small
\begin{tabular}{@{}r l c p{55mm}@{}}
  \toprule
  & Verbo & Requiere predicativo & Ejemplo \\
  \midrule
  1.  & \forma{parecer}   & sí & \respuesta{Pedro parece cansado.} \\
  2.  & \forma{escribir}  & & \\
  3.  & \forma{caer}      & & \\
  4.  & \forma{volverse}  & & \\
  5.  & \forma{construir} & & \\
  6.  & \forma{comer}     & & \\
  7.  & \forma{correr}    & & \\
  8.  & \forma{seguir}    & & \\
  9.  & \forma{resultar}  & & \\
  10. & \forma{haber}     & & \\
  11. & \forma{estar}     & & \\
  12. & \forma{ser}       & & \\
  13. & \forma{estudiar}  & & \\
  \bottomrule
\end{tabular}
\end{center}
\end{ejercicio}

\section{La sintaxis de los predicativos}\label{sec:09-sintaxis}

Como se presentó en el \capref{6}, en el marco de la gramática generativa, los
predicados sintácticos están compuestos por una proyección funcional que porta
la información flexiva (SF) y que contiene, a su vez, un sintagma verbal (SV),
que es el encargado de expresar la información eventiva.\footnote{Para más
detalles sobre las estructuras predicativas remitimos a
\textcite[211--250]{hernanz1987}.} Los predicativos, por su parte, son
semejantes a los predicados, pero carecen de toda esa articulación. Por eso, en
la bibliografía se los ha analizado en los términos de lo que se denomina
\termino{oración reducida} (OR), un tipo de estructura predicativa mínima que,
si bien expresa una predicación, carece de la manifestación de las categorías de
tiempo, aspecto y modo.\footnote{Para más detalles sobre las oraciones
reducidas, ver \textcite[425--429]{bosque2009}.} Dado que esa clase de análisis
requiere de una estructura más compleja, a los fines de este libro vamos a
marcar los predicativos simplemente del siguiente modo:

\begin{exe}
  \ex\label{ex:09-marcado}
    \begin{xlist}
      \ex Predicativo subjetivo obligatorio:\\
          {}[\etq{SF/O} [\etq{SN/Suj} Azucena Maizani] [\etq{SV/Pred}
          fue\etq{V/n} [\etq{SN/PSO} una\etq{Q/Esp} cantante\etq{N/n}]]]
      \ex Predicativo subjetivo no obligatorio:\\
          {}[\etq{SF/O} [\etq{SN/Suj} Carlos] [\etq{SV/Pred} volvió\etq{V/n}
          [\etq{SA/PSNO} cansado]]]
      \ex Predicativo objetivo obligatorio:\\
          {}[\etq{SF/O} [\etq{SN/Suj} la\etq{D/Esp} justicia\etq{N/n}]
          [\etq{SV/Pred} declaró\etq{V/n} [\etq{SN/COD} a\etq{MF} la\etq{D/Esp}
          acusada\etq{N/n}] [\etq{SA/POO} culpable]]]
      \ex Predicativo objetivo no obligatorio:\\
          {}[\etq{SF/O} [\etq{SN/Suj} Alfredo] [\etq{SV/Pred} compró\etq{V/n}
          [\etq{SN/COD} el\etq{D/Esp} disco\etq{N/n}] [\etq{SA/PONO} barato]]]
    \end{xlist}
\end{exe}

\begin{recuadro}[Para ampliar: análisis de oración reducida]
El análisis en términos de OR asume que el predicativo y el argumento al que se
aplica conforman una cláusula que no flexiona en tiempo, aspecto y modo. Para
obtener una oración completa, este análisis requiere algún mecanismo que
garantice que el sujeto de la oración reducida sea correferente con el sujeto u
objeto de la oración principal, según se trate de un predicativo subjetivo u
objetivo, respectivamente. Por ejemplo, en el caso del predicativo subjetivo
obligatorio, el sujeto se mueve desde la cláusula reducida, como se ve en el
siguiente árbol:

\begin{exe}
  \ex\label{ex:09-or1}
    \begin{xlist}
      \ex Azucena Maizani fue una cantante de tangos.
      \ex {}[\etq{SF/O} [\etq{SN/Suj} Azucena Maizani]\etq{i} [\etq{SV/Pred}
          fue\etq{V/n} [\etq{ORed/Comp} \sout{Azucena Maizani}\etq{i} [\etq{SN}
          una\etq{Q/Esp} cantante\etq{N/n} [\etq{SP/M} de\etq{P/n}
          [\etq{SN/Comp} tangos\etq{N/n}]]]]]]
      \ex
          \begin{forest} gramatica
            [SF / O
              [SN\etq{i}\\\textit{Azucena Maizani}]
              [SV
                [V\\\textit{fue}]
                [ORed
                  [SN\etq{i}\\{\color{gristexto}\textit{Azucena Maizani}}]
                  [SN [\textit{una cantante de tangos}, roof]]]]]
          \end{forest}
    \end{xlist}
\end{exe}

En este análisis, los dos elementos que aparecen en la oración reducida, la
huella de \forma{Azucena Maizani}, por un lado, y \forma{una cantante de
tangos}, por otro, no se etiquetan en relación con la función sintáctica que
cumplen, ya que la oración reducida carece de articulación interna. Los
predicativos subjetivos no obligatorios, como \exref{ex:09-or2}, se pueden
analizar también como oraciones reducidas.\pendiente{El original remite a (20)}
Sin embargo, a diferencia de lo que sucede con los obligatorios, en los no
obligatorios el argumento del cual se predica en la oración reducida no es un
elemento que luego se mueva, sino un pronombre nulo controlado, es decir, que es
obligatoriamente correferente con el sujeto de la oración principal. Llamaremos
a este pronombre PRO (no confundir con \textit{pro} en minúscula):

\begin{exe}
  \ex\label{ex:09-or2}
    \begin{xlist}
      \ex Carlos volvió cansado.
      \ex {}[\etq{SF/O} [\etq{SN/Suj} Carlos]\etq{i} [\etq{SV/Pred}
          volvió\etq{V/n} [\etq{ORed/Comp} PRO\etq{i} cansado\etq{SA}]]]
      \ex
          \begin{forest} gramatica
            [SF / O
              [SN\etq{i}\\\textit{Carlos}]
              [SV
                [V\\\textit{volvió}]
                [ORed
                  [PRO\etq{i}]
                  [SA\\\textit{cansado}]]]]
          \end{forest}
    \end{xlist}
\end{exe}

Análisis semejantes se pueden postular para los predicativos objetivos, tal como
se ilustra, respectivamente, en los siguientes ejemplos de predicativo objetivo
obligatorio y no obligatorio:

\begin{exe}
  \ex\label{ex:09-or3}
    \begin{xlist}
      \ex La justicia declaró a la acusada culpable.
      \ex {}[\etq{SF/O} [\etq{SN/Suj} la\etq{D/Esp} justicia\etq{N/n}]
          [\etq{SV/Pred} declaró\etq{V/n} [\etq{ORed/Comp} [\etq{SN/COD}
          a\etq{MF} la\etq{D/Esp} acusada\etq{N/n}] [\etq{SA} culpable]]]]
      \ex
          \begin{forest} gramatica
            [SF / O
              [SN [\textit{la justicia}, roof]]
              [SV
                [V\\\textit{declaró}]
                [ORed
                  [SN [\textit{a la acusada}, roof]]
                  [SA\\\textit{culpable}]]]]
          \end{forest}
    \end{xlist}
  \ex\label{ex:09-or4}
    \begin{xlist}
      \ex Alfredo compró el disco barato.
      \ex {}[\etq{SF/O} [\etq{SN/Suj} Alfredo] [\etq{SV/Pred} compró\etq{V/n}
          [\etq{SN/COD} el\etq{D/Esp} disco\etq{N/n}]\etq{i} [\etq{ORed}
          PRO\etq{i} barato\etq{SA}]]]
      \ex
          \begin{forest} gramatica
            [SF / O
              [SN\\\textit{Alfredo}]
              [SV
                [V\prima
                  [V\\\textit{compró}]
                  [SN\etq{i} [\textit{el disco}, roof]]]
                [ORed
                  [PRO\etq{i}]
                  [SA\\\textit{barato}]]]]
          \end{forest}
    \end{xlist}
\end{exe}
\end{recuadro}

\begin{ejercicio}\label{ej:09-analisis}
Haga el análisis sintáctico de las siguientes oraciones, según se muestra en el
ejemplo resuelto~1.

\begin{exe}
  \ex Ahora oficia de mecánico de bicicletas.
  \ex Mi hermano siempre toma el mate amargo.
  \ex Josefina desayuna callada.
  \ex Los niños escondieron culposos el jarrón roto.
  \ex En cada picadito de Argentina, un patadura se cree Messi.
  \ex Los violinistas callejeros se hacen los Paganini.
  \ex Me desperté muerto de hambre.
  \ex Las nuevas medidas dejaron desfinanciada la educación.
\end{exe}

\begin{resuelto}
1. [\etq{SF/O} [\etq{SN/Suj} \textit{pro}] [\etq{SV/Pred} [\etq{SAdv/ACT}
ahora\etq{Adv/n}] oficia\etq{V/n} [\etq{SP/PSO} de\etq{P/n} [\etq{SN/Comp}
mecánico\etq{N/n} [\etq{SP/M} de\etq{P/n} [\etq{SN/Comp}
bicicletas\etq{N/n}]]]]]]
\end{resuelto}
\end{ejercicio}

\section{Predicativos y homonimia estructural}\label{sec:09-homonimia}

Tal como se observó previamente, los predicativos no están limitados
categorialmente ---es decir, pueden ser sintagmas de cualquier clase--- y
establecen diferentes relaciones estructurales respecto del verbo que funciona
como núcleo del predicado y del argumento con el que se vinculan, sea el sujeto
o el objeto.\footnote{La homonimia estructural se introdujo en el \capref{2},
§\ref{sec:02-ambiguedad}, y en el \capref{5}, §\ref{sec:05-homonimia}.} Estas
dos características hacen que los predicativos tiendan a aparecer en oraciones
que son ambiguas por razones estructurales. El ejemplo de
\exref{ex:09-sentados} tiene tres lecturas posibles porque puede recibir tres
análisis sintácticos diferentes: en \exref{ex:09-sentados-a}, \forma{sentados}
forma parte del complemento objeto directo; en \exref{ex:09-sentados-b},
\forma{sentados} es un predicativo objetivo (los que están sentados son los
jugadores), y en \exref{ex:09-sentados-c} se trata de un predicativo subjetivo
(los que están sentados son los preparadores).

\begin{exe}
  \ex\label{ex:09-sentados} Los preparadores físicos observaban a los jugadores
      sentados.
    \begin{xlist}
      \ex\label{ex:09-sentados-a} {}[\etq{SF/O} [\etq{SN/Suj} los preparadores
          físicos] [\etq{SV/Pred} observaban [\etq{SN/COD} a los jugadores
          sentados]]]
      \ex\label{ex:09-sentados-b} {}[\etq{SF/O} [\etq{SN/Suj} los preparadores
          físicos] [\etq{SV/Pred} observaban [\etq{SN/COD} a los jugadores]
          [\etq{SA/PONO} sentados]]]
      \ex\label{ex:09-sentados-c} {}[\etq{SF/O} [\etq{SN/Suj} los preparadores
          físicos] [\etq{SV/Pred} observaban [\etq{SN/COD} a los jugadores]
          [\etq{SA/PSNO} sentados]]]
    \end{xlist}
\end{exe}
\pendiente{En los análisis, el original pone «Los entrenadores» como sujeto}

\begin{ejercicio}\label{ej:09-ambiguas}
Las siguientes oraciones son ambiguas. Explique mediante una paráfrasis cuáles
son las interpretaciones posibles y presente los análisis correspondientes para
cada una de esas lecturas, según se muestra en el ejemplo resuelto~1.

\begin{exe}
  \ex Trajeron a los perros enfermos.
  \ex Les entregaron los platos sucios.
  \ex Pedro pidió un postre borracho.
\end{exe}

\begin{resuelto}
1. a. [\etq{SF/O} [\etq{SN/Suj} \textit{pro}\etq{Pr/n}] [\etq{SV/Pred}
trajeron\etq{V/n} [\etq{SN/COD} a\etq{MF} los\etq{D/Esp} perros\etq{N/n}]
[\etq{SA/PO} enfermos\etq{A/n}]]]: \sig{Quienes estaban enfermos trajeron a los
perros}. b. [\etq{SF/O} [\etq{SN/Suj} \textit{pro}\etq{Pr/n}] [\etq{SV/Pred}
trajeron\etq{V/n} [\etq{SN/COD} a\etq{MF} los\etq{D/Esp} perros\etq{N/n}]
[\etq{SA/PSNO} enfermos\etq{A/n}]]]: \sig{Cuando trajeron a los perros, estos
estaban enfermos}. c. [\etq{SF/O} [\etq{SN/Suj} \textit{pro}\etq{Pr/n}]
[\etq{SV/Pred} trajeron\etq{V/n} [\etq{SN/COD} a\etq{MF} los\etq{D/Esp}
perros\etq{N/n} enfermos\etq{A/M}]]]: \sig{Trajeron a los perros, que estaban
enfermos}.\pendiente{Las etiquetas de a y b están intercambiadas respecto de las
paráfrasis (a es PSNO; b, PONO); el original rotula las lecturas a, b, b}
\end{resuelto}
\end{ejercicio}

\begin{ejercicio}\label{ej:09-vinas}
Lea atentamente el siguiente fragmento de \textit{Un dios cotidiano}, de
D.~Viñas. Luego identifique todos los predicativos y señale de qué tipo de
predicativo se trata en cada caso.

% Fragmento literario protegido: se incorpora a mano desde el original.
\textoaincorporar{el fragmento de \textit{Un dios cotidiano}, de D.~Viñas, desde
«Que-lo-nios, dije tres o cuatro veces» hasta «\ldots{}la voz me salía pastosa,
la matizaba»}
\end{ejercicio}

\begin{recuadro}[Reflexión final]
En este capítulo hemos abordado un tipo particular de función del sintagma
verbal: el predicativo. Al respecto, hemos observado que los predicativos pueden
corresponder a diferentes categorías: SN, SA, SP e incluso SAdv. Hemos
explorado también la tipología básica de los predicativos, que se desarrolla en
torno a dos parámetros: i) si están orientados al sujeto o al objeto y ii) si
son obligatorios o no obligatorios. Esto lleva a establecer cuatro tipos de
predicativos: subjetivo obligatorio (PSO), subjetivo no obligatorio (PSNO),
objetivo obligatorio (POO) y objetivo no obligatorio (PONO). Hemos también
presentado cómo se los analiza sintácticamente y qué propiedades sintácticas
tiene cada subclase.
\end{recuadro}

% --- Secciones finales -------------------------------------------------
\lecturasrecomendadas

\solucionario

\begin{solucion}{\ref{ej:09-predicados}}
2. Eduarda Mansilla [\etq{SV} \subrayado{escribió varios libros}]. 3. Eduarda
Mansilla [\etq{SV} \subrayado{ofició de traductora}]. 4. Eduarda Mansilla [fue
también [\etq{SN} \subrayado{compositora}]]. 5. Margarita Abella Caprile
[\etq{SV} \subrayado{nació en 1901}]. 6. La poesía de Margarita Caprile [era
[\etq{SA} \subrayado{muy íntima}]]. 7. Margarita Caprile [fue [\etq{SN}
\subrayado{amiga de Alfonsina Storni}]]. 8. María Luisa Carnelli [era [\etq{SP}
\subrayado{de La Plata}]]. 9. María Luisa Carnelli [\etq{SV}
\subrayado{firmaba sus letras de tango con seudónimo}].
\end{solucion}

\begin{solucion}{\ref{ej:09-primaria}}
3. Predicación primaria: \forma{La sala está llena}. Predicación secundaria: no
hay. 4. Predicación primaria: \forma{El niño lloró}. Predicación secundaria:
\forma{El niño estaba desconsolado}. 5. Predicación primaria: \forma{Las
Malvinas son argentinas}. Predicación secundaria: no hay. 6. Predicación
primaria: \forma{Ella respondió la llamada}. Predicación secundaria: \forma{Ella
estaba muy tranquila}. 7. Predicación primaria: \forma{Martín se siente solo}.
Predicación secundaria: no hay.
\end{solucion}

\begin{solucion}{\ref{ej:09-tabla}}
a y b. 1a. [\etq{Suj} Santos Vega] fue [\etq{SN} un payador argentino]. 1b.
[\etq{Suj} Los demás payadores] consideraban [\etq{COD} a Santos Vega]
[\etq{SA} invencible]. 2a. [\etq{Suj} Gabino Ezeiza] murió [\etq{SA} pobre] en
el barrio de Flores. 2b. Se suele considerar [\etq{COD} a José Betinotti]
[\etq{SN} el nexo entre los payadores y los primeros tangueros]. 3a. Para Gabino
Ezeiza, [\etq{Suj} Pancho Luna] ofició [\etq{SP} de maestro]. 3b. [\etq{Suj}
José Betinotti] convirtió [\etq{SP} en exitosa] [\etq{COD} la línea de las
payadas que exaltaban el amor materno]. 4a. [\etq{Suj} Arturo de Nava] solía
vestir [\etq{SA} elegante]. 4b. [\etq{Suj} El Mercosur] declaró [\etq{SN}
Patrimonio Cultural] [\etq{COD} al arte de la payada]. 5a. [\etq{Suj} José
Betinotti] siempre recordaba [\etq{SA} agradecido] [\etq{COD} a su amigo Pablo
Vázquez] en cada función. 5b. [\etq{Suj} El alcalde] acusó [\etq{COD} a Juan
Moreira] [\etq{SP} de mentiroso].

c. La marca morfosintáctica que visibiliza la relación entre el predicativo y el
sujeto o el objeto es la concordancia en género y número, en los casos en que el
núcleo del predicativo es un adjetivo. En los casos en que el predicativo es un
SN o un SP, la relación que se establece es de compatibilidad de rasgos
semánticos, pero no una concordancia propiamente dicha. En el contexto de los
SSPP, esa compatibilidad se da con el complemento de la preposición.
\end{solucion}

\begin{solucion}{\ref{ej:09-ps-po}}
3. José Betinotti se declaraba [\etq{PS} admirador de Alfredo Palacios]. 4. Para
una grabación particular, Mariano Villar Sáenz Peña convocó [\etq{PO} de
guitarrista] a Rosita Quiroga. 5. Las sucesivas crisis económicas dejaron al
Café de los Angelitos [\etq{PO} muy deteriorado]. 6. Ángel Vargas tenía un
pulmón [\etq{PO} muy comprometido]. 7. Alfredo Palacios y Juan B.~Justo solían
discutir [\etq{PS} enfervorizados] en el Café de los Angelitos. 8. Rosita Quiroga
se retiró de su carrera artística [\etq{PS} todavía joven]. 9. En sus últimos
años de vida, a Enrique Santos Discépolo se lo veía [\etq{PS}
solo].\pendiente{Con \textit{se lo veía}, \textit{solo} se predica del objeto
(\textit{lo})} 10. Azucena Maizani murió [\etq{PS} casi olvidada] en 1970.
\end{solucion}

\begin{solucion}{\ref{ej:09-pso-psno}}
3. Ernesto es [\etq{PSO} biólogo]. 4. Carlos salió [\etq{PSNO} despeinado] de su
casa. 5. Mi perro corre [\etq{PSNO} medio rengo] cuando se asusta. 6. Ese reloj
parece [\etq{PSO} nuevo]. 7. Rocío miró la foto [\etq{PSNO} asombrada]. 8.
Enrique escribió su novela [\etq{PSNO} muy entusiasmado]. 9. Martina resultó
[\etq{PSO} la ganadora del concurso]. 10. La actriz interpretó el papel
[\etq{PSNO} muy segura].
\end{solucion}

\begin{solucion}{\ref{ej:09-sonono}}
i) Hay predicativos solo en 1b, 2b, 2d, 3b, 4b, 4c, 4e, 5b, 6b, 7b y 8a.
\pendiente{El original omite 1b de la lista, aunque lo resuelve, y analiza 2b como
predicativo en iii)}

ii) 2d. Sigue [\etq{SA/PO} molesto] por lo que le dije. 3b. Se puso [\etq{SA/PO}
rojo]. 4b. Volvió [\etq{SA/PNO} enfermo]. 4c. Volvió [\etq{SA/POO} loco] a su
amigo. 4e. Se volvió [\etq{SA/PO} loco]. 5b. Juan se sintió [\etq{SA/PO}
cansado]. 6b. Esteban se hizo [\etq{SN/PO} el tonto]. 7b. Ese tipo me cae
[\etq{SA/PSO} simpático]. 8a. Jugó [\etq{SA/PNO} muy concentrado].

iii) 2a. Sigue [\etq{SN/COD} su instinto]. 2b. Sigue [\etq{SAdv/PO} mal]. 2c.
Sigue [\etq{SP/ACL} en Madrid]. 3a. Se puso [\etq{SN/COD} un sombrero]. 4a.
Volvió [\etq{SAdv/ACT} temprano]. 4d. Volvió [\etq{SN/ACT} una noche]. 5a. Juan
se sintió [\etq{SN/COD} olor a humo] después de hacer el asado. 6a. Esteban se
hizo [\etq{SN/COD} el desayuno]. 7a. Ese tipo cayó [\etq{SAdv/ACM} mal]. Por eso
se rompió [\etq{SN/COD} el tobillo]. 8b. Jugó [\etq{SN/COD} el partido].
\end{solucion}

\begin{solucion}{\ref{ej:09-pruebas}}
2. Prueba de reemplazo por \forma{lo} invariable. 3. Prueba de coordinación
entre la predicación primaria y la secundaria. 4. Prueba de concordancia. 5.
Prueba de movimiento: misma explicación que en el ejemplo 1.
\end{solucion}

\begin{solucion}{\ref{ej:09-enriquez}}
[\etq{CCL} cerca del puerto] / [\etq{PSO} australiano] / [\etq{PSNO} de mecánico
de bicicletas] / [\etq{ApE} capital de Australia Occidental] / [\etq{M} a larga
distancia] / [\etq{PSO} insólitamente estable, permanente y feliz] / [\etq{PSO}
la capital de estado más aislada del mundo] / [\etq{PSO} la ciudad de más de un
millón de habitantes más aislada del mundo] / [\etq{PSO} comprobable y cierto] /
[\etq{PSO} una localidad muy cool, algo estrafalaria, favorita de mochileros y
\textit{newagers}] / [\etq{M} con jardines llenos de cristales curativos y
hadas].
\end{solucion}

\begin{solucion}{\ref{ej:09-significado}}
2. i) a y c: [\etq{PSO} estupefacta]; b no presenta ningún predicativo. ii) En
el primer y el tercer uso, \forma{permanecer} refiere a que \sig{la gente no
salió de su asombro}; en el segundo caso, se dice que \sig{la gente no se
trasladó espacialmente}. 3. i) a: [\etq{PSO} aislada]; b no presenta ningún
predicativo. ii) En a, el predicativo describe el estado de Perth; en b, se
indica su ubicación espacial. 4. i) a no presenta ningún predicativo; b:
[\etq{PO} culpable]. ii) En a, el verbo \forma{declarar} significa \sig{hizo una
declaración}; en b, la predicación indica que \sig{el juez condenó al imputado}.
\end{solucion}

\begin{solucion}{\ref{ej:09-vf}}
ii. F. iii. F. iv. V.
\end{solucion}

\begin{solucion}{\ref{ej:09-verbos}}
a. No requieren predicativo: \forma{escribir, caer, construir, comer, correr,
haber, estudiar}; requieren predicativo: \forma{volverse, seguir, resultar,
estar, ser}. b. Respuestas posibles: 4. El profesor se volvió loco. 8. Mi
hermano sigue en cama / enfermo. 9. La clase de gramática resultó muy
interesante. 11. La sopa está caliente. 12. Comer frutas por la mañana es
verdaderamente importante.
\end{solucion}

\begin{solucion}{\ref{ej:09-analisis}}
2. [\etq{SF/O} [\etq{SN/Suj} mi\etq{D/Esp} hermano\etq{N/n}] [\etq{SV/Pred}
[\etq{SAdv/ACT} siempre\etq{Adv/n}] toma\etq{V/n} [\etq{SN/COD} el\etq{D/Esp}
mate\etq{N/n}] [\etq{SA/PONO} amargo\etq{A/n}]]].

3. [\etq{SF/O} [\etq{SN/Suj} Josefina\etq{N/n}] [\etq{SV/Pred} desayuna\etq{V/n}
[\etq{SA/PSNO} callada\etq{A/n}]]].

4. [\etq{SC=SF/O} [\etq{SN/Suj} los\etq{D/Esp} niños\etq{N/n}] [\etq{SV/Pred}
escondieron\etq{V/n} [\etq{SA/PSNO} culposos\etq{A/n}] [\etq{SN/COD}
el\etq{D/Esp} jarrón\etq{N/n} roto\etq{A/M}]]].

5. [\etq{SF/O} [\etq{SP/ACL} en\etq{P/n} [\etq{SN/Comp} cada\etq{Q/Esp}
picadito\etq{N/n} [\etq{SP/M} de\etq{P/n} [Argentina\etq{N/n}]]]]\etq{i},
[\etq{SN/Suj} un\etq{Q/Esp} patadura\etq{N/n}] [\etq{SV/Pred} se\etq{IPCR}
cree\etq{V/n} [\etq{SN/PSO} Messi\etq{N/n}] h\etq{i}]]. \textit{Aclaración}: el
constituyente \forma{En cada picadito de Argentina} se halla desplazado de la
posición en la que se origina, interna al predicado.

6. [\etq{SF/O} [\etq{SN/Suj} los\etq{D/Esp} violinistas\etq{N/n}
callejeros\etq{A/M}] [\etq{SV/Pred} se\etq{IPCR} hacen\etq{V/n} [\etq{SN/PSO}
los\etq{D/Esp} Paganini\etq{N/n}]]]. \textit{Aclaración}: \forma{Paganini}
aparece recategorizado como nombre común.

7. [\etq{SF/O} [\etq{SN/Suj} \textit{pro}] [\etq{SV/Pred} me\etq{IPCR}
desperté\etq{V/n} [\etq{SA/PSNO} muerto\etq{A/n} [\etq{SP/M} de\etq{P/n}
[\etq{SN/Comp} hambre\etq{N/n}]]]]].

8. [\etq{SF/O} [\etq{SN/Suj} las\etq{D/Esp} nuevas\etq{A/M} medidas\etq{N/n}]
[\etq{SV/Pred} dejaron\etq{V/n} [\etq{SA/POO} desfinanciada\etq{A/n}]
[\etq{SN/COD} la\etq{D/Esp} educación\etq{N/n}]]].
\end{solucion}

\begin{solucion}{\ref{ej:09-ambiguas}}
En cada caso incluimos una paráfrasis de modelo. Eso no implica que no haya
otras paráfrasis posibles.

2. a. [\etq{SF/O} [\etq{SN/Suj} \textit{pro}] [\etq{SV/Pred} les\etq{Pr/COI}
entregaron\etq{V/n} [\etq{SN/COD} los\etq{D/Esp} platos\etq{N/n}] [\etq{SA/PO}
sucios\etq{A/n}]]]: \sig{Les entregaron sucios los platos}. b. [\etq{SF/O}
[\etq{SN/Suj} \textit{pro}] [\etq{SV/Pred} les\etq{Pr/COI} entregaron\etq{V/n}
[\etq{SN/COD} los\etq{D/Esp} platos\etq{N/n} sucios\etq{A/M}]]]: \sig{Les
entregaron los platos que estaban sucios}.

3. a. [\etq{SF/O} [\etq{SN/Suj} Pedro\etq{N/n}] [\etq{SV/Pred} pidió\etq{V/n}
[\etq{SN/COD} un\etq{D/Esp} postre\etq{N/n}] [\etq{SA/PSNO}
borracho\etq{A/n}]]]: \sig{Pedro pidió el postre estando borracho}. b.
[\etq{SF/O} [\etq{SN/Suj} Pedro\etq{N/n}] [\etq{SV/Pred} pidió\etq{V/n}
[\etq{SN/COD} un\etq{D/Esp} postre borracho\etq{N/n}]]]:\footnote{\forma{Postre
borracho} en 3b funciona como un sustantivo compuesto.} \sig{Pedro pidió un
postre llamado «postre borracho»}.
\end{solucion}

\begin{solucion}{\ref{ej:09-vinas}}
Era una palabra [\etq{PSO} divertida]. / La ballena me salió [\etq{PSO}
estupenda]. / Parecía [\etq{PSO} el comienzo de una firma ampulosa, una «Q»
inflada, potente]. / El ornitorrinco provocó asombro: era [\etq{PSO} un bicho
ridículo y conmovedor]. / Sentía la mano [\etq{POO} ligera]. / Yo era [\etq{PSO}
eficaz], sabía, podía dibujar\ldots{} la voz me salía [\etq{PSO} pastosa], la
matizaba.
\end{solucion}
